\section{LLM Backend System: DeepSeek API, Claude CLI Subprocess, and Backend-Aware Task Generation}
\label{sec:backend}

\subsection{Dual-Backend Architecture}

UMAF's LLM layer implements a dual-backend architecture providing transparent access to two fundamentally different invocation paradigms: a traditional API-based approach via DeepSeek's OpenAI-compatible endpoint, and a subprocess-based approach that shells out to the Claude Code CLI. The architecture unifies both backends behind a common \texttt{LLMProvider} ABC with a single \texttt{invoke(messages) $\rightarrow$ AIMessage} interface.

\subsection{DeepSeekProvider}

\raggedright
The DeepSeek backend \cite{deepseekv3} wraps LangChain's \texttt{ChatOpenAI} pointed at \path{https://api.deepseek.com/v1} with model \texttt{deepseek-chat}, temperature 0.3, and \texttt{max\_tokens=8192}. API key resolution uses \texttt{DEEPSEEK\_API\_KEY} from \texttt{.env} via \texttt{python-dotenv}. The \texttt{invoke()} method delegates directly to LangChain's \texttt{ChatOpenAI.invoke()}, which handles OpenAI-compatible request/response serialization, retry, and error mapping.

\subsection{ClaudeCLILLM}

The Claude CLI backend \cite{mcp} constructs a \texttt{claude -p <prompt>} subprocess with \texttt{--output-format stream-json}, \texttt{--verbose}, and \texttt{--permission-mode bypassPermissions}. Key features include:
\begin{itemize}
    \item \textbf{Timeout}: Default 600s, enforced via \texttt{threading.Timer} that kills the subprocess on timeout
    \item \textbf{Environment injection}: Merges 12 env vars from \texttt{claude\_env\_sample.json}
    \item \textbf{Working directory}: Accepts \texttt{cwd} parameter for path resolution
    \item \textbf{Allowed tools}: Per-role tool restriction via \texttt{--allowedTools}
    \item \textbf{Two invocation modes}: \texttt{invoke()} for text output, \texttt{stream\_invoke()} for stream-json event processing
\end{itemize}

\subsection{Backend-Aware Task Generation}

\raggedright
This is UMAF's most important architectural innovation. Every \texttt{AgentRole.build\_task(backend, ...)} receives the backend identifier and tailors the prompt accordingly at three layers:

\textbf{Layer 1 --- BaseAgent Prompt Builders}: DeepSeek receives a 311-line prompt with JSON tool-call formatting rules, 4 escape rules, and a write strategy hierarchy (prefer \texttt{run\_command} Python one-liners over \texttt{write\_file}). Claude CLI receives a 30-line prompt with native tool instructions.

\textbf{Layer 2 --- AgentRole Per-Backend Branching}: For DeepSeek: emphasizes \texttt{call\_claude} for deep reasoning, uses \texttt{download\_file} + \texttt{read\_file} for arxiv.org. For Claude CLI: the agent IS a Claude Code instance and must NOT spawn nested \texttt{claude -p} calls.

\textbf{Layer 3 --- BaseDecomposerRole Instructions}: Head agents append backend-specific suffixes. Claude CLI agents are told ``Use your own knowledge---do NOT search the web.'' DeepSeek agents receive web search instructions and arxiv.org pre-fetch patterns.

\subsection{The Pre-Fetch Layer}

Claude Code's cc-switch blocks domain verification for arxiv.org. UMAF's pre-fetch layer downloads academic content at the framework level via Python \texttt{urllib} before agents run. The mechanism: (1) search for subtask via \texttt{web\_search()}, (2) extract arxiv.org URLs via regex, (3) download up to 3 via \texttt{download\_file()} to local HTML, (4) inject file paths into Claude CLI worker prompts with ``Read them directly'' instructions.

\subsection{CheckpointManager and Conversation Logger}

\texttt{CheckpointManager} provides versioned state persistence at \path{agent_log/{name}_v{version:02d}_checkpoint.json}. The \texttt{load\_previous(n)} method enables cross-version context reuse for retry---loading the highest version below $n$ with full message history restoration. The conversation logger writes timestamped JSON logs for debugging, performance analysis, and self-evolution feedback.

\subsection{Backend Trade-offs}

\begin{table}[H]
\centering
\small
\caption{DeepSeek vs.\ Claude CLI Comparison}
\label{tab:backend}
\begin{tabular}{@{}p{3cm}p{5.2cm}p{5.2cm}@{}}
\toprule
\textbf{Dimension} & \textbf{DeepSeek (ChatOpenAI)} & \textbf{Claude CLI (Subprocess)} \\
\midrule
Tool Calling & JSON parsing from text (4-strategy parser) & Native tool calling via stream-json \\
Reliability & Moderate---prone to JSON malformation & High---structurally validated by CLI \\
Multi-Step Tasks & Gets stuck in loops, loses state & Maintains coherent state across contexts \\
Cost & \textasciitilde\$0.28/M input, \textasciitilde\$0.42/M output & \textasciitilde\$3/M input, \textasciitilde\$15/M output \\
Latency & Fast---single API call per iteration & Slower---subprocess startup overhead \\
Circuit Breakers & Full intervention system & Limited to hard timeout + auto-retry \\
Checkpointing & After every tool execution & After every stream-json event \\
\bottomrule
\end{tabular}
\end{table}

\subsection{JSON Tool-Call Parsing Pipeline}

Since DeepSeek lacks native tool calling \cite{deepseekv3,deepseekr1}, UMAF implements a 4-strategy parser:
\begin{enumerate}
    \item Markdown code fences (reversed order---newest first)
    \item Standard JSON key order (\texttt{\{"tool": ..., "args": ...\}})
    \item Reversed key order (\texttt{\{"args": ..., "tool": ...\}})
    \item JSON repair: remove trailing commas, escape raw whitespace via state machine
\end{enumerate}
A write strategy hierarchy mitigates JSON escaping: prefer \texttt{run\_command} with Python one-liner, then \texttt{write\_lines}, then \texttt{write\_file} as last resort.
